\documentclass[a4paper,12pt]{article}

% ===========================================================================
%  Лабораторная работа 3.1.1 «Магнитометр»
%
%  Сборка:  cd 3.1.1/report && latexmk -pdf main.tex
%  Таблицы читаются из ../data/processed/*.csv, а итоговые числа --
%  из ../data/processed/results.tex; оба готовит ../scripts/process_data.py.
%  Пока данных нет, на месте чисел стоит ???, на месте текста -- «?? НЕОБХОДИМО ЗАПОЛНИТЬ ??».
% ===========================================================================

% ============================================================
% Общий преамбул для отчётов по лабораторным работам (labs_lpr)
% Подключается из report/main.tex командой:
%   % ============================================================
% Общий преамбул для отчётов по лабораторным работам (labs_lpr)
% Подключается из report/main.tex командой:
%   % ============================================================
% Общий преамбул для отчётов по лабораторным работам (labs_lpr)
% Подключается из report/main.tex командой:
%   \input{../../common/latex_template/preamble.tex}
% ============================================================

% Поддержка русского языка и кодировки
\usepackage[T2A]{fontenc}
\usepackage{lmodern}
\usepackage[utf8]{inputenc}
\usepackage[russian]{babel}

% Математические пакеты
\usepackage{amsmath, amssymb}

% Настройка полей
\usepackage{geometry}
\geometry{top=2cm, bottom=2cm, left=2.5cm, right=1.5cm}

% Колонтитулы
\usepackage{fancyhdr}
\usepackage{graphicx}

% Современные таблицы
\usepackage{tabularray}

% Работа с csv файлами
\usepackage{csvsimple-l3}
\usepackage{tikz}

% Графики
\usepackage{pgfplots}
\pgfplotsset{compat=1.18}

% Ссылки
% \usepackage{hyperref} % гипперссылки, нужно оформление вместо красного квадрата
\usepackage[russian]{cleveref} % умные ссылки, нужно переоформление для формул, рисунков и тд (вместо ф-л что-то нужно придумать)

% ============================================================

% Поддержка русского языка и кодировки
\usepackage[T2A]{fontenc}
\usepackage{lmodern}
\usepackage[utf8]{inputenc}
\usepackage[russian]{babel}

% Математические пакеты
\usepackage{amsmath, amssymb}

% Настройка полей
\usepackage{geometry}
\geometry{top=2cm, bottom=2cm, left=2.5cm, right=1.5cm}

% Колонтитулы
\usepackage{fancyhdr}
\usepackage{graphicx}

% Современные таблицы
\usepackage{tabularray}

% Работа с csv файлами
\usepackage{csvsimple-l3}
\usepackage{tikz}

% Графики
\usepackage{pgfplots}
\pgfplotsset{compat=1.18}

% Ссылки
% \usepackage{hyperref} % гипперссылки, нужно оформление вместо красного квадрата
\usepackage[russian]{cleveref} % умные ссылки, нужно переоформление для формул, рисунков и тд (вместо ф-л что-то нужно придумать)

% ============================================================

% Поддержка русского языка и кодировки
\usepackage[T2A]{fontenc}
\usepackage{lmodern}
\usepackage[utf8]{inputenc}
\usepackage[russian]{babel}

% Математические пакеты
\usepackage{amsmath, amssymb}

% Настройка полей
\usepackage{geometry}
\geometry{top=2cm, bottom=2cm, left=2.5cm, right=1.5cm}

% Колонтитулы
\usepackage{fancyhdr}
\usepackage{graphicx}

% Современные таблицы
\usepackage{tabularray}

% Работа с csv файлами
\usepackage{csvsimple-l3}
\usepackage{tikz}

% Графики
\usepackage{pgfplots}
\pgfplotsset{compat=1.18}

% Ссылки
% \usepackage{hyperref} % гипперссылки, нужно оформление вместо красного квадрата
\usepackage[russian]{cleveref} % умные ссылки, нужно переоформление для формул, рисунков и тд (вместо ф-л что-то нужно придумать)


% Пакеты, нужные именно этой работе
\usepackage{subcaption}
\usepackage{circuitikz}
\usetikzlibrary{arrows.meta, calc, angles, quotes}

% --- Служебные макросы ------------------------------------------------------
% Короткий пропуск (число, которое надо вписать после измерений)
\newcommand{\qq}{\textbf{???}}
% Пропуск, требующий осмысленного текста
\newcommand{\fillme}{\textbf{?? НЕОБХОДИМО ЗАПОЛНИТЬ ??}}
% Имя файла моноширинным шрифтом (символы «_» и т.п. экранировать не нужно)
\newcommand{\filename}[1]{\texttt{\detokenize{#1}}}
% Рисунок из report/img/; если файла ещё нет -- рамка-заглушка
\newcommand{\labimage}[2]{%
    \IfFileExists{img/#1}{%
        \includegraphics[width=#2\linewidth]{img/#1}%
    }{%
        \fbox{\parbox[c][0.4\dimexpr#2\linewidth\relax][c]{%
                \dimexpr#2\linewidth-2\fboxsep-2\fboxrule\relax}{%
            \centering\small\filename{img/#1}\\[2mm]\fillme}}%
    }%
}
% Итоговые числа из ../data/processed/results.tex (пишет process_data.py):
%   \res{key} -> «значение \pm погрешность», \resv{key} -> значение,
%   \ress{key} -> погрешность, \rese{key} -> относительная погрешность, %.
% Ключи -- см. data/processed/fit_results.csv. Нет числа -- выводится ???.
\InputIfFileExists{../data/processed/results.tex}{}{}
\newcommand{\resget}[1]{\ifcsname res@#1\endcsname\csname res@#1\endcsname\else\qq\fi}
\newcommand{\res}[1]{\resget{#1}}
\newcommand{\resv}[1]{\resget{#1@v}}
\newcommand{\ress}[1]{\resget{#1@s}}
\newcommand{\rese}[1]{\resget{#1@e}}

% Колонтитулы (как в остальных работах репозитория)
\pagestyle{fancy}
\fancyhf{}
\fancyhead[L]{%
    \IfFileExists{../../common/latex_template/miptlogo.png}{%
        \includegraphics[height=1cm]{../../common/latex_template/miptlogo.png}%
    }{%
        \textbf{[ЛОГОТИП МФТИ]}%
    }%
}
\fancyhead[R]{\textbf{Емельянов ДА, Байдин ФМ}}
\fancyfoot[C]{\thepage}
\renewcommand{\headrulewidth}{0.4pt}
\setlength{\headheight}{33pt}

\begin{document}

% --- ТИТУЛЬНЫЙ ЛИСТ ---
\begin{titlepage}
\begin{center}
    {\large МОСКОВСКИЙ ФИЗИКО-ТЕХНИЧЕСКИЙ ИНСТИТУТ}
\end{center}
\begin{center}
    {\large Физтех-школа физики и исследований им. Ландау\break ЛФИ}
\end{center}

    \vspace{3cm}

\vspace{0.15cm}
{
\begin{center}
    \huge \bf Лабораторная работа №3.1.1 по курсу Общая физика\\
    \vspace{0.4cm}
    \LARGE Магнитометр.
\end{center}
}
\vspace{1.5cm}
\begin{flushright}
Емельянов Дмитрий Алексеевич\\
Байдин Фаддей Михайлович\\
\vspace{0.4cm}
Б02-502
\end{flushright}
\vspace{10cm}
\begin{center}
    \large \qq{} 2026, г. Долгопрудный
\end{center}
\end{titlepage}
% ----------------------

\section{Аннотация}
Цель работы --- определить горизонтальную составляющую магнитного поля Земли
$B_0$ и установить количественное соотношение между единицами
электрического тока в системах СИ и СГС, т.\,е. измерить
электродинамическую постоянную $c$.

Поле $B_0$ находится с помощью электромагнитного магнитометра: измеряются
угол отклонения магнитной стрелки в поле намагниченного стержня и период
крутильных колебаний того же стержня в поле Земли, что позволяет исключить
неизвестный магнитный момент стержня. Затем через витки магнитометра
пропускается ток периодической разрядки конденсатора: по отклонению стрелки
ток определяется в единицах СИ (по магнитному действию), а по
протекающему заряду --- в гауссовых единицах. Отношение этих величин даёт
$c$. Ожидается, что $B_0$ по порядку величины совпадёт со справочным
значением для Московского региона (с поправкой на искажения поля в здании),
а $c$ --- со скоростью света в вакууме.

\section{Теоретические сведения}

\subsection{Поле стержня и кругового тока}
Магнитометр (\cref{fig:magnetometer}) состоит из $N$ последовательно
соединённых круговых витков радиуса $R$, расположенных вертикально; в
центре кольца на тонкой нити подвешена короткая магнитная стрелка. В
отсутствие других полей стрелка устанавливается вдоль горизонтальной
составляющей поля Земли $B_0$ (в плоскости магнитного меридиана), и кольцо
поворотом совмещают с этой плоскостью.

Поле намагниченного стержня с магнитным моментом $\mathfrak{m}$ вдали от
него приближённо совпадает с полем точечного диполя
\begin{equation}
    \label{eq:dipole}
    \mathbf{B}(\mathbf{r}) = \frac{\mu_0}{4\pi}
    \left(\frac{3(\boldsymbol{\mathfrak{m}}\cdot\mathbf{r})\mathbf{r}}{r^5}
          - \frac{\boldsymbol{\mathfrak{m}}}{r^3}\right).
\end{equation}
Стержень вставляется в гнездо Р на горизонтальном диаметре кольца так, что
центр кольца лежит на оси, перпендикулярной стержню, на расстоянии $R$:
\begin{equation}
    \label{eq:B1}
    B_1 = \frac{\mu_0}{4\pi}\frac{\mathfrak{m}}{R^3}.
\end{equation}
Поле в центре кольца с током $I$ по закону Био--Савара
\begin{equation}
    \label{eq:B2}
    B_2 = \frac{\mu_0 I}{2R}N.
\end{equation}
Оба поля перпендикулярны плоскости кольца, т.\,е. плоскости меридиана, и
в области стрелки однородны, так как размеры стрелки много меньше $R$.

\subsection{Отклонение стрелки}
При появлении перпендикулярного меридиану поля $B_\perp$ ($B_1$ или $B_2$)
стрелка устанавливается вдоль суммарного поля и поворачивается на угол
$\varphi$:
\begin{equation}
    \label{eq:tan}
    B_\perp = B_0\operatorname{tg}\varphi.
\end{equation}
Угол измеряется по смещению $x$ светового зайчика, отражённого от
зеркальца на стрелке, на шкале, удалённой на расстояние $L$
(\cref{fig:magnetometer}). Отражённый луч поворачивается на $2\varphi$, и при
малых углах
\begin{equation}
    \label{eq:tan_x}
    \operatorname{tg}\varphi \approx \frac{x}{2L}.
\end{equation}

\subsection{Горизонтальная составляющая поля Земли}
Чтобы исключить неизвестный момент $\mathfrak{m}$, измеряется период малых
крутильных колебаний стержня, подвешенного за середину в поле $B_0$
(упругостью нити пренебрегаем). Уравнение колебаний
$J\ddot\alpha + \mathfrak{m}B_0\alpha = 0$ даёт
\begin{equation}
    \label{eq:period}
    T = 2\pi\sqrt{\frac{J}{\mathfrak{m}B_0}},
\end{equation}
где момент инерции цилиндрического стержня массы $m$, длины $l$ и радиуса
$r$ относительно оси, проходящей через середину перпендикулярно стержню,
\begin{equation}
    \label{eq:inertia}
    J = m\left(\frac{l^2}{12} + \frac{r^2}{4}\right)
      = \frac{ml^2}{12}\left[1 + 3\left(\frac{r}{l}\right)^2\right].
\end{equation}
Из \cref{eq:B1,eq:tan,eq:tan_x,eq:period} при $x = x_1$
\begin{equation}
    \label{eq:B0}
    B_0 = \frac{1}{TR}\sqrt{\frac{2\pi\mu_0 J L}{R\,x_1}}
    \quad[\text{СИ}].
\end{equation}

\subsection{Ток в единицах СИ и СГС}
Пропуская ток через витки и измеряя отклонение $x_2$, из
\cref{eq:B2,eq:tan,eq:tan_x} находим ток в единицах СИ (это соответствует
определению ампера через магнитное действие тока):
\begin{equation}
    \label{eq:I_SI}
    I = A\operatorname{tg}\varphi_2 = \frac{2B_0R}{\mu_0N}\,\frac{x_2}{2L}
    \quad[\text{А}],
    \qquad A = \frac{2B_0R}{\mu_0N},
\end{equation}
где $A$ --- постоянная прибора в данном месте комнаты.

Тот же ток можно измерить по переносимому заряду (определение единицы тока
в гауссовой системе). Если конденсатор ёмкости $C$, заряженный до
напряжения $U$, $\nu$ раз в секунду разряжать через витки, средний ток
\begin{equation}
    \label{eq:I_CGS}
    I = CU\nu \quad[\text{ед. СГС}],
\end{equation}
где $C$ выражается в сантиметрах ($1~\text{Ф} \approx 9\cdot10^{11}$~см), а
$U$ --- в единицах СГС ($1~\text{В} \approx \frac{1}{300}$~ед. СГС).
Отношение численных значений одного и того же тока даёт
электродинамическую постоянную
\begin{equation}
    \label{eq:c}
    c\left[\frac{\text{м}}{\text{с}}\right]
    = \frac{1}{10}\,\frac{I[\text{СГС}]}{I[\text{СИ}]}.
\end{equation}
Соотношение \cref{eq:c} следует из определений единиц: $1~\text{А} =
\frac{c[\text{см/с}]}{10}$~ед. СГС тока.

\subsection{Основные допущения}
\begin{itemize}
    \item стержень --- точечный диполь (его длина мала по сравнению с $R$);
    \item поля $B_1$, $B_2$ однородны в области стрелки;
    \item поле в месте измерения периода $T$ совпадает с полем в центре
          магнитометра (измерения проводятся рядом с прибором);
    \item упругость нити подвеса стержня и стрелки пренебрежимо мала;
    \item углы отклонения малы, $\operatorname{tg}2\varphi\approx2\operatorname{tg}\varphi$;
    \item за время между замыканиями вибратора конденсатор успевает полностью
          зарядиться и разрядиться.
\end{itemize}

\section{Методика измерений}

\subsection{Схема установки}
Схема магнитометра и оптической системы показана на
\cref{fig:magnetometer}, фотография установки --- на \cref{fig:photo}. Осветитель О освещает два зеркала:
З$_1$ на стрелке (подвижный зайчик) и З$_2$, жёстко связанное с кольцом
(неподвижный зайчик). Совмещение зайчиков поворотом кольца означает, что
плоскость витков совпадает с плоскостью магнитного меридиана.

\begin{figure}[!ht]
    \centering
    \begin{tikzpicture}[scale=1.0, >=Stealth, font=\small]
        % кольцо (вид сверху -- отрезок в плоскости меридиана)
        \draw[line width=2.2pt, gray] (-2.6,0) -- (2.6,0);
        \node[above left] at (-2.0,0.05) {К};
        % гнездо Р со стержнем
        \fill[black] (2.45,-0.35) rectangle (2.75,0.35);
        \node[right] at (2.8,0) {Р (стержень)};
        % стрелка с зеркалом
        \draw[line width=1.4pt, rotate=25] (-0.6,0) -- (0.6,0);
        \draw[line width=1.2pt, blue!70!black, rotate=25] (-0.05,-0.3) -- (-0.05,0.3);
        \node[below right] at (0.15,-0.05) {С, З$_1$};
        % поля
        \draw[->, thick] (0,0) -- (2.0,0) node[below] {$B_0$};
        \draw[->, thick] (0,0) -- (0,0.93) node[left] {$B_\perp$};
        \draw[->, thick, red!70!black] (0,0) -- (2.0,0.93) node[above right] {$B_\Sigma$};
        \draw (0.9,0) arc (0:25:0.9);
        \node at (1.1,0.22) {$\varphi$};
        % шкала и осветитель
        \draw[line width=1.6pt] (-2.6,-4) -- (2.6,-4);
        \foreach \x in {-2.5,-2.25,...,2.5} \draw (\x,-4) -- (\x,-3.9);
        \node[below] at (-2.4,-4.05) {Ш};
        \node[draw, minimum width=8mm, minimum height=5mm] (O) at (-0.9,-4.6) {О};
        % лучи
        \draw[dashed, orange!80!black] (0,-4) -- (0,0);
        \draw[->, orange!80!black, thick] (0,0) -- (1.86,-4);
        \draw (0,-1.2) arc (-90:-65:1.2);
        \node at (0.28,-1.55) {$2\varphi$};
        \draw[<->] (0,-4.35) -- (1.86,-4.35) node[midway, below] {$x$};
        \draw[<->] (-3.0,0) -- (-3.0,-4) node[midway, left] {$L$};
    \end{tikzpicture}
    \caption{Схема магнитометра (вид сверху) и измерения угла отклонения
             стрелки: К --- витки, С --- стрелка с зеркалом З$_1$,
             Р --- гнездо для стержня, Ш --- шкала, О --- осветитель}
    \label{fig:magnetometer}
\end{figure}

\begin{figure}[!ht]
    \centering
    \labimage{setup_photo.jpg}{0.6}
    \caption{Фотография установки (по желанию)}
    \label{fig:photo}
\end{figure}

Для измерения тока витки подключаются к схеме \cref{fig:circuit}:
электромагнитное реле (вибратор) с частотой $\nu$ попеременно подключает
конденсатор $C$ то к источнику (зарядка до напряжения $U$, измеряемого
вольтметром), то к виткам магнитометра (разрядка).

\begin{figure}[!ht]
    \centering
    \begin{circuitikz}[font=\small]
        \draw (0,0) to[V, l=$\mathcal{E}$] (0,3) -- (3.4,3);
        \draw (3.4,3) node[ocirc]{};
        \draw (4.6,3) node[ocirc]{} -- (5.6,3)
              to[nos, l=ключ] (7.4,3)
              to[L, l=витки К] (7.4,0) -- (0,0);
        \draw (4,2.4) node[circ]{} -- (3.48,2.93);
        \node[above] at (4,3.1) {реле, $\nu$};
        \draw (4,2.4) to[C, l_=$C$] (4,0);
        \draw (4,2.4) -- (2.4,2.4) to[voltmeter, l_=$U$] (2.4,0);
    \end{circuitikz}
    \caption{Схема питания витков магнитометра}
    \label{fig:circuit}
\end{figure}

\subsection{Порядок измерений}
\begin{enumerate}
    \item Совмещаются зайчики (кольцо --- в плоскости меридиана).
    \item В гнездо Р вставляется стержень, измеряется смещение $x_1$; стержень
          переворачивается, измеряется смещение в другую сторону. Расхождение
          не должно превышать 5\,\%.
    \item Измеряется расстояние $L$ от шкалы до зеркала.
    \item Рядом с магнитометром измеряется время $t$ нескольких ($N_\text{кол}$)
          крутильных колебаний стержня, $T = t/N_\text{кол}$. Число
          колебаний выбирается так, чтобы $\sigma_T/T \le 1\,\%$:
          $N_\text{кол} \ge \sigma_t/(0{,}01\,T)$.
    \item Штангенциркулем измеряются длина $l$ и диаметр $d = 2r$ стержня.
    \item По \cref{eq:inertia,eq:B0} рассчитывается $B_0$.
    \item Собирается схема \cref{fig:circuit}; при работающем вибраторе
          измеряются $U$ и отклонение $x_2$ при двух полярностях.
    \item По \cref{eq:I_SI,eq:I_CGS,eq:c} вычисляются токи и $c$.
\end{enumerate}

Обработка выполняется скриптом \filename{scripts/process_data.py}.
Погрешности косвенных величин находятся по формулам для степенных
зависимостей; например, для \cref{eq:B0}
\begin{equation}
    \label{eq:sigma_B0}
    \varepsilon_{B_0} = \sqrt{\varepsilon_T^2 + \left(\tfrac32\varepsilon_R\right)^2
        + \left(\tfrac12\varepsilon_J\right)^2 + \left(\tfrac12\varepsilon_L\right)^2
        + \left(\tfrac12\varepsilon_{x_1}\right)^2},
    \qquad
    \varepsilon_J = \sqrt{\varepsilon_m^2 + (2\varepsilon_l)^2},
\end{equation}
где $\varepsilon_a = \sigma_a/a$. Погрешность прямых многократных измерений
складывается из случайной (СКО среднего) и приборной. Поскольку $I[\text{СИ}]$
вычисляется через $B_0$, при оценке $\varepsilon_c$ учитывается, что после
подстановки \cref{eq:B0} в \cref{eq:I_SI}
$I[\text{СИ}] \propto T^{-1}R^{-1/2}J^{1/2}L^{-1/2}x_1^{-1/2}x_2$.

\section{Используемое оборудование}

\begin{table}[!ht]
    \centering
    \caption{Измерительные приборы и их погрешности}
    \label{tab:equipment}
    \begin{tblr}{colspec={|X[2,l]|X[2,l]|X[1.4,c]|}, hlines, vlines,
                 row{1}={font=\bfseries}, width=\textwidth}
        Прибор & Измеряемая величина & Погрешность \\
        Магнитометр (осветитель, шкала) & смещение зайчика $x$ & $\sigma_x = \qq$~мм \\
        Рулетка & расстояние $L$ & $\sigma_L = \qq$~см \\
        Штангенциркуль & $l$, $d$ стержня & $\qq$~мм \\
        Секундомер & время колебаний $t$ & $\sigma_t = \qq$~с \\
        Вольтметр \qq & напряжение $U$ & $\sigma_U = \qq$~В \\
        Конденсатор (на установке) & ёмкость $C = \qq$~мкФ & $\qq$ \\
        Электромагнитный переключатель (вибратор) & частота $\nu = \qq$~Гц & $\qq$ \\
        Стержень (на установке) & масса $m = \qq$~г & $\qq$ \\
    \end{tblr}
\end{table}

Также используются источник питания ($U \approx 90\div100$~В), ключ и
прибор для определения периода крутильных колебаний (стеклянный сосуд со
спицей и нитью). Параметры магнитометра: $R = \qq$~см, $N = \qq$.

Предварительная оценка: наибольший вклад в погрешность $B_0$ ожидается от
радиуса кольца (входит в степени $3/2$) и периода $T$; в погрешность $c$ ---
от $T$, $x_2$ и ёмкости $C$. \fillme

\section{Результаты измерений и обработка данных}

\subsection{Параметры установки}
Все параметры установки записаны в \filename{data/raw/setup_params.csv}:
\begin{itemize}
    \item радиус кольца $R = \qq$~см, число витков $N = \qq$;
    \item расстояние от шкалы до зеркала $L = \qq$~см;
    \item масса стержня $m = \qq$~г;
    \item ёмкость $C = \qq$~мкФ, частота вибратора $\nu = \qq$~Гц.
\end{itemize}

\subsection{А. Горизонтальная составляющая магнитного поля Земли}

\paragraph{Отклонение стрелки в поле стержня.}
Результаты измерения смещения зайчика при двух ориентациях стержня
приведены в \cref{tab:deflection}.

\begin{table}[!ht]
    \centering
    \caption{Смещение подвижного зайчика в поле намагниченного стержня;
             $\delta$ --- относительное расхождение $x_1^+$ и $x_1^-$}
    \label{tab:deflection}
    \csvreader[
        tabularray = {
            colspec = {|c|c|c|c|c|},
            hlines, vlines,
            row{1} = {font=\bfseries},
        },
        table head = {
            \textbf{№} & \textbf{$x_1^+$, мм} & \textbf{$x_1^-$, мм}
            & \textbf{$x_1$, мм} & \textbf{$\delta$, \%} \\
        }
    ]%
    {../data/processed/deflection_rod_processed.csv}%
    {trial=\csvN, x1_plus_mm=\csvXp, x1_minus_mm=\csvXm, x1_mm=\csvX, dev_percent=\csvDev}%
    {%
        ${}\csvN$ & ${}\csvXp$ & ${}\csvXm$ & ${}\csvX$ & ${}\csvDev$
    }
\end{table}

Среднее смещение $x_1 = \res{xone}$~мм ($\varepsilon = \rese{xone}$\,\%),
$\operatorname{tg}\varphi_1 = x_1/2L = \res{tgone}$.

\paragraph{Момент инерции стержня.}
Размеры стержня приведены в \cref{tab:rod}.

\begin{table}[!ht]
    \centering
    \caption{Размеры намагниченного стержня}
    \label{tab:rod}
    \csvreader[
        tabularray = {
            colspec = {|c|c|c|},
            hlines, vlines,
            row{1} = {font=\bfseries},
        },
        table head = {
            \textbf{№} & \textbf{$l$, мм} & \textbf{$d$, мм} \\
        }
    ]%
    {../data/processed/rod_size_processed.csv}%
    {trial=\csvN, l_mm=\csvL, d_mm=\csvD}%
    {%
        ${}\csvN$ & ${}\csvL$ & ${}\csvD$
    }
\end{table}

$l = \res{lrod}$~мм, $d = \res{drod}$~мм. Поправка $3(r/l)^2 = \res{corr}$\,\%.
По \cref{eq:inertia}
\begin{equation}
    J = (\res{J})\cdot10^{-7}~\text{кг}\cdot\text{м}^2,
    \qquad \varepsilon_J = \rese{J}\,\%.
\end{equation}

\paragraph{Период крутильных колебаний.}
Результаты приведены в \cref{tab:period}.

\begin{table}[!ht]
    \centering
    \caption{Период крутильных колебаний стержня в поле Земли;
             $\sigma_T = \sigma_t/N_\text{кол}$ --- приборная погрешность}
    \label{tab:period}
    \csvreader[
        tabularray = {
            colspec = {|c|c|c|c|c|},
            hlines, vlines,
            row{1} = {font=\bfseries},
        },
        table head = {
            \textbf{№} & \textbf{$N_\text{кол}$} & \textbf{$t$, с}
            & \textbf{$T$, с} & \textbf{$\sigma_T$, с} \\
        }
    ]%
    {../data/processed/rod_period_processed.csv}%
    {trial=\csvN, N_osc=\csvNosc, t_s=\csvT, T_s=\csvPer, sigma_T_s=\csvSPer}%
    {%
        ${}\csvN$ & ${}\csvNosc$ & ${}\csvT$ & ${}\csvPer$ & ${}\csvSPer$
    }
\end{table}

\begin{equation}
    T = \res{T}~\text{с}, \qquad \varepsilon_T = \rese{T}\,\%.
\end{equation}

\paragraph{Результат.}
По \cref{eq:B0,eq:sigma_B0}
\begin{equation}
    \label{eq:B0_result}
    B_0 = \res{Bzero}~\text{мкТл} = \res{Bzerogs}~\text{Гс},
    \qquad \varepsilon_{B_0} = \rese{Bzero}\,\%.
\end{equation}
Магнитный момент стержня, найденный из \cref{eq:period},
$\mathfrak{m} = \res{mrod}$~мА$\cdot$м$^2$.

Справочное значение горизонтальной составляющей поля Земли для
Долгопрудного (модель WMM или IGRF): $B_0^\text{спр} = \qq$~мкТл.

\subsection{Б. Электродинамическая постоянная}
Постоянная прибора \cref{eq:I_SI}: $A = \res{A}$~мА. Результаты измерения
тока приведены в \cref{tab:current}; $x_2$ --- среднее отклонение при двух
полярностях.

\begin{table}[!ht]
    \centering
    \caption{Измерение тока через витки магнитометра двумя способами}
    \label{tab:current}
    \small
    \csvreader[
        tabularray = {
            colspec = {|c|c|c|c|c|c|c|c|c|},
            hlines, vlines,
            row{1} = {font=\bfseries},
        },
        table head = {
            \textbf{№} & \textbf{$U$, В} & \textbf{$x_2^+$, мм} & \textbf{$x_2^-$, мм}
            & \textbf{$x_2$, мм} & \textbf{$\operatorname{tg}\varphi_2$}
            & \textbf{$I_\text{СИ}$, мА} & \textbf{$I_\text{СГС}$, $10^6$ ед.}
            & \textbf{$c$, $10^8$ м/с} \\
        }
    ]%
    {../data/processed/current_processed.csv}%
    {trial=\csvN, U_V=\csvU, x2_plus_mm=\csvXp, x2_minus_mm=\csvXm, x2_mm=\csvX,
     tg_phi2=\csvTg, I_SI_mA=\csvIsi, I_CGS_1e6=\csvIcgs, c_1e8=\csvC}%
    {%
        ${}\csvN$ & ${}\csvU$ & ${}\csvXp$ & ${}\csvXm$ & ${}\csvX$ & ${}\csvTg$
        & ${}\csvIsi$ & ${}\csvIcgs$ & ${}\csvC$
    }
\end{table}

По средним значениям $U = \res{U}$~В и $x_2 = \res{xtwo}$~мм:
\begin{equation}
    I[\text{СИ}] = \res{ISI}~\text{мА},
    \qquad
    I[\text{СГС}] = (\res{ICGS})\cdot10^{6}~\text{ед. СГС}.
\end{equation}
По \cref{eq:c}
\begin{equation}
    \label{eq:c_result}
    c = (\res{c})\cdot10^8~\text{м/с},
    \qquad \varepsilon_c = \rese{c}\,\%.
\end{equation}
Отклонение от табличного значения $c = 2{,}998\cdot10^8$~м/с составляет
$\res{cdev}$\,\%.

\section{Обсуждение результатов}
\begin{itemize}
    \item Точность $B_0$: основной вклад в $\varepsilon_{B_0}$ вносит \fillme.
    \item Сравнение $B_0 = \res{Bzero}$~мкТл со справочным значением
          $\qq$~мкТл: \fillme. Магнитометр находится в железобетонном здании,
          поэтому измеряемое поле может заметно отличаться от поля Земли
          (арматура, токонесущие провода, приборы), а также различаться в
          месте измерения $T$ и в центре магнитометра.
    \item Значение $c$ не зависит от того, совпадает ли $B_0$ с полем Земли:
          важно лишь, чтобы при измерении $x_1$, $T$ и $x_2$ поле было одним
          и тем же. \fillme
    \item Применимость допущений: длина стержня $l = \res{lrod}$~мм по
          сравнению с $R$ (дипольное приближение); малость углов
          $\varphi_1$, $\varphi_2$; полная разрядка конденсатора за время
          $1/(2\nu)$. \fillme
    \item Возможные источники ошибок и способы их уменьшения: \fillme.
\end{itemize}

\section{Заключение}
\begin{itemize}
    \item Горизонтальная составляющая магнитного поля в месте установки:
          $B_0 = \res{Bzero}$~мкТл ($\varepsilon = \rese{Bzero}$\,\%);
          справочное значение $\qq$~мкТл.
    \item Электродинамическая постоянная:
          $c = (\res{c})\cdot10^8$~м/с ($\varepsilon = \rese{c}$\,\%),
          что \fillme{} с табличным значением.
\end{itemize}
\fillme

\appendix
\section{Лабораторный журнал}
\labimage{journal.jpg}{1.0}

\end{document}
